% ===================================================================
%  TABELL Nr. 4 — Mannesalter an Alter.
%  Traduction 2026 d'apres texto/io, controlee sur texto/fr, texto/de
%  et texto/nl. Consigne : voir texto/lb/10-tabell-01.tex.
%
%  LE SAC DE CHASSE EN BANDOULIERE EST ENTIEREMENT EN GRAS, comme
%  dans les vingt et une colonnes qui l'ont deja corrige :
%  « Juegdtäsch iwwer der Schëller ».
%
%  LA PARENTE EST LE PLUS PUR MORCEAU DE LUXEMBOURGEOIS DU LIVRET,
%  et il tombe justement dans le tableau que la colonne galicienne
%  avait trouve le plus proche de la langue voisine. Ni l'allemand ni
%  le francais ne disent ce que dit cette page :
%
%    « Monni »      l'oncle, « Tatta » la tante ;
%    « Bomi »       la grand-mere, « Bopa » le grand-pere ;
%    « Peiter »     le parrain, « Giedel » la marraine ;
%    « Schwéiesch » la belle-soeur ;
%    « Weesekand »  l'orphelin.
%
%  Ces mots ne sont ni des emprunts ni des composes savants : ce sont
%  les mots que les enfants apprennent avant de savoir lire, et c'est
%  precisement pour cela qu'ils ont resiste aux deux voisins. LE
%  CRITERE DU LIEU D'APPRENTISSAGE DIT ICI SON CAS LE PLUS SIMPLE :
%  ce qu'on apprend a la maison ne s'emprunte pas.
%
%  ET LA MALADIE DIT LE CONTRAIRE, DANS LA MEME PAGE. « Rheumatismus »,
%  « Longenentzündung », « Ordonnance », « Mixtur », « Apdikt »,
%  « Dokter » viennent du dehors ; seuls « Houscht » la toux,
%  « Féiwer » la fievre, « Zännwéi » le mal de dents et
%  « Erkältung » le rhume sont a lui. C'est le partage que la colonne
%  lituanienne avait trouve au meme tableau — le mal qu'on sent et le
%  mal qu'on diagnostique — et il se verifie ici dans une langue
%  d'une autre famille.
% ===================================================================

%%K t04-apar-1 apar
\VUsaut{4.0mm}
\VUtitre{15.0pt}{60}{TABELL Nr. 4}
\VUsaut{1.6mm}
\VUtitre{11.4pt}{0}{Mannesalter an Alter.}
\VUsaut{3.2mm}

%%K t04-tit-1 sub
\VUcentre{11.4pt}{0}{\textit{Éischt Zeen.}}
\VUsaut{1.6mm}
\VUtitre{11.4pt}{0}{D'Hochzäitsiessen.}
\VUsaut{2.6mm}

%%K t04-tit-2 sub
\VUpk{10.2pt}{40}{Den Hierscht.}
\VUsaut{3.0mm}

%%K t04-c1-01-1 p
1. --- Déi Jonk sinn grouss ginn. Ee vun hinnen, de Jang, bestued sech. Mir sinn
beim \VUgras{Hochzäitsiessen} dobäi, dat ronderëm dee selwechten Dësch d'Famillje
vun de Bestuedene versammelt.

%%K t04-c1-02-1 p
2. --- Mir sinn am \VUgras{Hierscht,} am \VUgras{Mount September.} De kale Wand vum
\VUgras{Oktober} a vum \VUgras{November} huet dem Land säi grénge Blieder-Kleed nach
net geraf. D'Owesloft ass lau a gebeemt; dofir ass d'Iessen ënner der
\VUgras{Veranda} \textsuperscript{(8)} nieft dem Gaart.

%%K t04-c1-03-1 p
3. --- Wäit ewech, um \VUgras{Hiwwel} \textsuperscript{(3)}, steet d'Heem vun de
Elteren vum Jang, en elegant \VUgras{Schlass} \textsuperscript{(1)}, dat am Gréngen
verstoppt ass; schéi Wéngerten, déi ganz gudde Wäi ginn, bedecken den Hang vum
Hiwwelchen. De Wënzer baut se un a këmmert sech ëm si.

%%K t04-c1-04-1 p
4. --- Iwwert de Wéngerte fléit e Schwarm \VUgras{Rebhénger} \textsuperscript{(2)},
déi vum Nokomme vun der \VUgras{Juegdopsiicht} \textsuperscript{(5)} opgeschreckt goufen.
Dëse leeft mat senger \VUgras{Flënt} ënnert dem Aarm an der \VUgras{Juegdtäsch iwwer
der Schëller} d'\VUgras{Gestrëpps} \textsuperscript{(4)} mat sengem \VUgras{Hond}
\textsuperscript{(6)} of. Hie passt op d'Gebitt op a verhënnert, datt d'Wilddéif
d'Wëld ëmbréngen.

%%K t04-c1-05-1 p
5. --- Baussen op der \VUgras{Veranda} klëmmt e \VUgras{Wäistack-Espalier,} vun deem
dicht \VUgras{Wäintrauwen} \textsuperscript{(7)} erofhänken.

%%K t04-c1-06-1 p
6. --- Déi schéin Hierschtblummen, déi den \VUgras{Dësch} \textsuperscript{(16)}
schmécken, sinn \VUgras{Chrysanthemen} \textsuperscript{(9)}; de \VUgras{Papp}
\textsuperscript{(10)} vun der Braut huet eng dervun a säi Knappfluch gestach.

%%K t04-c1-07-1 p
7. --- D'Iessen geet op en Enn. De \VUgras{Bedéngten} \textsuperscript{(11)} fänkt un,
d'Dessertteller ze bréngen, a wäert geschwënn déi verschidde Stécker vum Gedeck
ewechhuelen --- \VUgras{Läffelen} \textsuperscript{(14)}, \VUgras{Forschetten}
\textsuperscript{(12)}, \VUgras{Messeren} \textsuperscript{(13)},
\VUgras{Plateauen} \textsuperscript{(15)} ---, déi ee net méi brauch.

%%K t04-tit-3 sub
\VUpk{10.2pt}{40}{D'Verwandtschaft.}
\VUsaut{3.0mm}

%%K t04-c1-08-1 p
8. --- All Leit gesinn e frou aus, an d'Laachen, mat deem d'\VUgras{Mamm}
\textsuperscript{(17)} vum Bräitchemann hir Kanner ukuckt, weist hire Gléck.

%%K t04-c1-09-1 p
9. --- Dës Hochzäit verbënnt zwou räich Famillje vum Land. De Jang, de
\VUgras{Bräitchemann} \textsuperscript{(18)}, ass de \VUgras{Jong} vun engem groussen
Grondbesëtzer, an hie gëtt de \VUgras{Schwéiersoun} vun engem geschéckten
Industriellen.

%%K t04-c1-10-1 p
10. --- D'\VUgras{Bestuedene} si jonk, an trotzdeem waren si scho laang
\VUgras{verlobt.} Déi \VUgras{jonk Fra} \textsuperscript{(19)}, d'Martha, ass
charmant an hirem \VUgras{wäisse Kleed,} dat mat Orangebléie geschmückt ass.

%%K t04-c1-11-1 p
11. --- Hire \VUgras{Schwéierpapp} \textsuperscript{(20)} ass ganz glécklech, esou eng
frëndlech \VUgras{Schwéierduechter} ze hunn, an d'\VUgras{Schwéiermamm}
\textsuperscript{(21)} vum Jang laacht, well si hir \VUgras{Duechter} esou schéi
gesäit.

%%K t04-c1-12-1 p
12. --- Um Enn vum Dësch sëtzen déi aner \VUgras{Invitéen} \textsuperscript{(22)}:
\VUgras{Verwandt, Frënn, Cousinen, Cousinnen.} E \VUgras{Maître d'hôtel}
\textsuperscript{(23)}, mat der Serviette ënnert dem Aarm, bedéngt si fläisseg.

%%K t04-tit-4 sub
\VUpk{10.2pt}{40}{D'Frënn.}
\VUsaut{3.0mm}

%%K t04-c1-13-1 p
13. --- Riets vum Dësch, do ass e \VUgras{Meedchen} \textsuperscript{(24)}, eng
\VUgras{Frëndin} vun der jonker Braut, duerno hire \VUgras{Brudder,} deen hiren
\VUgras{Éierebegleeder} \textsuperscript{(25)} ass. Hie schwätzt mat senger
\VUgras{Éierendamm} \textsuperscript{(26)}, der Schwëster vum Bräitchemann, déi duerch
dës Hochzäit d'\VUgras{Schwéiesch} vun der Martha gëtt. Si ass eng charmant jonk
Fra. Si huet schéi \VUgras{Bijouen, eng Pärelekette} an en gëllent \VUgras{Armband.}

%%K t04-c1-14-1 p
14. --- Nieft hinnen ass den Här, deen opsteet a säi Glas hieft, e \VUgras{Frënd}
\textsuperscript{(27)} vun der Famill. Hien ass ee vun de \VUgras{Zeie vum
Bräitchemann.} Hie \VUgras{bréngt en Toast aus,} dréckt op d'\VUgras{Gesondheet} vun
de jonke Bestuedenen a wënscht hinne laangt Gléck. Hien ass festlech ugedoen, mat
engem Frack a mat \VUgras{Lackschong} \textsuperscript{(28)}. Hien ass de Begleeder
vun der \VUgras{Bomi} \textsuperscript{(29)}, déi \VUgras{Witfra} ass.

%%K t04-c1-15-1 p
15. --- De jonke Mann am \VUgras{Frack} \textsuperscript{(31)}, mat engem dënnen
schwaarze Schnauzert, ass de \VUgras{Schwoer} \textsuperscript{(30)} vum Jang. Hien
ass en hervirroungend Ingenieur. Hien ass den Noper vun enger ganz eleganter
\VUgras{jonker Damm} \textsuperscript{(33)}, der \VUgras{eelerer Schwëster} vun der
Martha an der Mamm vum zierleche klenge Meedchen, dat kënnt, um Aarm vu sengem
Cousin, fir eng Blumm ze bréngen an him e \VUgras{gudden Owend ze wënschen.} Dës Damm
huet ganz schéin \VUgras{Ouerréng} an eng elegant \VUgras{Kappbedeckung.}

%%K t04-c1-16-1 p
16. --- De klenge Cousin, esou komesch mat senger \VUgras{Bluse}
\textsuperscript{(36)} a senger baucheger \VUgras{Pofferbox} \textsuperscript{(37)},
ass ganz bedauerlech. Hien ass e \VUgras{Weesekand} \textsuperscript{(35)}. Ganz jonk
huet hie säi Papp a seng Mamm verluer. Säi Monni ass säi \VUgras{Vormond}
\textsuperscript{(38)}, an hien ass onbestuet bliwwen, fir dat aarmt Kand
z'erzéien. Hien ass e gudde Mann. Hien ass e bësse kal a ganz kuerzsiichteg; hie
benotzt e \VUgras{(Nues-)Pince-nez.}

%%K t04-tit-5 sub
\VUsaut{2.4mm}
\VUpk{10.2pt}{40}{Den Dessert.}
\VUsaut{3.0mm}

%%K t04-c1-17-1 p
17. --- Hannert de Gäscht, op engem Dësch mat engem \VUgras{Dëschduch} bedeckt, ass
den \VUgras{Dessert} \textsuperscript{(39)} virbereet. An enger grousser Schuel, do
sinn Uebst: \VUgras{Pëschingen} \textsuperscript{(40)}, \VUgras{Bieren}
\textsuperscript{(41)}, \VUgras{Äppel} \textsuperscript{(42)},
\VUgras{Aprikosen} \textsuperscript{(43)} a \VUgras{Wäintrauwen}
\textsuperscript{(44)}. Nieft drun \VUgras{Ananas} \textsuperscript{(45)}, e schéine
\VUgras{Kuch} \textsuperscript{(46)}, \VUgras{Likörfläschen} \textsuperscript{(48)} a
\VUgras{Champagner} \textsuperscript{(49)}, an och Sailen vu propere
\VUgras{Telleren} \textsuperscript{(47)}.

%%K t04-c1-18-1 p
18. --- De \VUgras{Kellermeeschter} \textsuperscript{(50)} mécht eng
\VUgras{Fläsch} \textsuperscript{(52)} alen Wäin mam \VUgras{Kuerkenzéier}
\textsuperscript{(51)} op. De \VUgras{Kuerken} \textsuperscript{(53)} schéngt ganz
schwéier erauszekréien. Dëse Kellermeeschter huet eng \VUgras{Serviette}
\textsuperscript{(54)} ënnert dem Aarm.

%%K t04-c1-19-1 p
19. --- Nom Iessen ginn d'Hären an d'Fëmmerzëmmer, an d'Damme ginn sech fir de Bal
prett maachen.

%%K t04-tit-6 sub
\VUsaut{5.0mm}
\VUcentre{11.4pt}{0}{\textit{Zweet Zeen.}}
\VUsaut{1.6mm}
\VUpk{11.4pt}{40}{De Gebuertsdag vum Bopa.}
\VUsaut{2.6mm}

%%K t04-tit-7 sub
\VUpk{10.2pt}{40}{De Besuch.}
\VUsaut{3.0mm}

%%K t04-c2-01-1 p
1. --- Zéng Joer si vergaangen, dem Jang seng Elteren si al ginn a hunn hir dräi
Enkelkanner ganz gär, zwee \VUgras{Enkelen} \textsuperscript{(2)} an eng
\VUgras{Enkelin} \textsuperscript{(3)}. Well haut de \VUgras{Gebuertsdag vum Bopa}
ass, kommen d'Kanner hinne gratuléieren.

%%K t04-c2-02-1 p
2. --- De klenge Péiter ass nach ganz jonk, hien ass eréischt dräi Joer al. Hie
gesäit d'\VUgras{Kaz} \textsuperscript{(1)} noukommen. Hie wëll hire schéinen
säidegen \VUgras{Pelz} an hire laange \VUgras{Schwanz} sträichelen; hien denkt net
drun, datt ënnert deene zierleche \VUgras{Patten} béis spëtz Krallen verstoppt sinn,
déi hie vläicht krazen.

%%K t04-c2-03-1 p
3. --- Seng Schwëster d'Gabriell, déi him d'Hand gëtt, ass vill eeschter, an de
Marcel mat sengem groussen \VUgras{Mantel} \textsuperscript{(4)} a sengem
liedernen \VUgras{Gürtel} \textsuperscript{(5)} gesäit e bësse männlech aus, trotz
senger kurzer Box, déi seng \VUgras{Strëmp} \textsuperscript{(6)} gesi léisst.

%%K t04-c2-04-1 p
4. --- Hire Papp huet säin décken \VUgras{Wanteriwwerzéier} \textsuperscript{(7)} mat
dem \VUgras{Pelzkolli} \textsuperscript{(8)} ugedoen, an an senger \VUgras{Täsch}
\textsuperscript{(9)} huet hien e Fichu. Seng Fra huet hire Filzhutt, mat
\VUgras{Fiederen} \textsuperscript{(10)} geschmückt, hire \VUgras{Jackett}
\textsuperscript{(11)}, hire \VUgras{Pelzboa} \textsuperscript{(12)} an hire
\VUgras{Muff} \textsuperscript{(13)}.

%%K t04-c2-05-1 p
5. --- Et ass ganz kal, well de Wanter herrscht. De \VUgras{Schnéi}
\textsuperscript{(19)} ass de ganzen Dag gefall, an d'Diecher vun den Haiser si ganz
wäiss. Zesumme mat der \VUgras{Nuecht} gëtt d'Keelt nach méi schaarf. Um Himmel
blénken de \VUgras{Mound} \textsuperscript{(17)} an d'\VUgras{Stären}
\textsuperscript{(18)} an duerchdréngen d'\VUgras{Däischtert.}

%%K t04-tit-8 sub
\VUsaut{4.2mm}
\VUpk{10.2pt}{40}{D'Krankheet.}
\VUsaut{3.0mm}

%%K t04-c2-06-1 p
6. --- Deen aarme \VUgras{Blannen} \textsuperscript{(20)}, deen iwwert de Plaz
virun der \VUgras{Apdikt} \textsuperscript{(16)} geet, gefouert vu sengem
\VUgras{Pudel} \textsuperscript{(21)}, ass ganz onglécklech. D'Justine,
d'\VUgras{Déngschtmeedchen} \textsuperscript{(14)}, bréngt aus der Kichen eng
\VUgras{Taass} \textsuperscript{(15)} ganz waarme \VUgras{Kraiderdéi,} deen
generéis gezockert ass.

%%K t04-c2-07-1 p
7. --- D'\VUgras{Bomi} \textsuperscript{(23)}, déi hiren \VUgras{Ouerbrëll}
\textsuperscript{(26)} um Dësch sichen wëll, fir d'\VUgras{Rezept}
\textsuperscript{(27)} ze liesen, dat den \VUgras{Dokter} \textsuperscript{(25)} grad
geschriwwen huet, steet vun hirem Fauteuil op an stäipt sech op hire \VUgras{Stack}
\textsuperscript{(24)}.

%%K t04-c2-08-1 p
8. --- Si leidt dacks u \VUgras{klengen Onwuelsinn, Schwindel, Erkältungen,
Zännwéi;} hir Been si schwaach, an hir Aen sinn net méi ganz gutt. Trotzdeem mécht si
sech gär lëschteg iwwer hiren alen \VUgras{Dokter,} dee hir ëmmer eng
\VUgras{Mixtur} \textsuperscript{(28)} verschreiwe wëll. Haut ass si ganz frou, well
si hir jonk Famill ronderëm sech huet.

%%K t04-c2-09-1 p
9. --- Et ass gemittlech an der gutt zouener Stuff. Am marmorne \VUgras{Kamäin}
\textsuperscript{(29)} flackert e \VUgras{schéint Feier} \textsuperscript{(30)}, an ee
fillt sech glécklech am mëllen \VUgras{Liicht} vun der \VUgras{Lamp}
\textsuperscript{(31)}, déi ee grad ugefaangen huet.

%%K t04-tit-9 sub
\VUsaut{4.2mm}
\VUpk{10.2pt}{40}{De Bopa.}
\VUsaut{3.0mm}

%%K t04-c2-10-1 p
10. --- Souguer de \VUgras{Bopa} \textsuperscript{(33)} huet vergiess, datt hie krank
war, datt säi \VUgras{Rheumatismus} hien un säi \VUgras{Fauteuil} \textsuperscript{(34)}
festhält, an datt hien nach gëschter \VUgras{Féiwer a Schwächeaanfäll} hat. Hien
erënnert sech net méi drun, datt hien de leschte Wanter un enger
\VUgras{Longenentzündung} gelidden huet, an datt e rauen a schmäerzhaften
\VUgras{Houscht} hie ganz gequält huet.

%%K t04-c2-10-2 p
An trotzdeem ass hie ganz schwéier gesond ginn. Vun elo un geet et him e bësse
besser. Mä säi Been, dat hien op e \VUgras{Këssen} \textsuperscript{(38)} stäipt, dat
op engem klenge \VUgras{Schemel} \textsuperscript{(39)} läit, deet him och nach wéi.

%%K t04-c2-11-1 p
11. --- Wann hie suergfälteg an säi waarme \VUgras{Schlofrock} \textsuperscript{(35)}
mat décke \VUgras{Knäpp} \textsuperscript{(36)} aus Pärelmutt agewéckelt ass, a wann
hien nieft sech dat kleng \VUgras{Weesemeedchen} \textsuperscript{(40)} huet, dat him
d'Zeitung virliest, an dat mat enger wäisser \VUgras{Kap,} engem bloen
\VUgras{Kleed} \textsuperscript{(42)} an enger \VUgras{Pelerine}
\textsuperscript{(41)} ugedoen ass, da beschwéiert hie sech net.

%%K t04-c2-12-1 p
12. --- All Joer, wann de \VUgras{Kalenner} \textsuperscript{(43)} nees den
\VUgras{Datum} vum éischten \VUgras{Dezember} weist, waart hien ongedëlleg op seng
\VUgras{Enkelkanner,} well hie weess, datt si dësen wichtegen \VUgras{Datum} net
vergiessen.

%%K t04-c2-13-1 p
13. --- Hien erënnert sech gerührt un déi ganz al Zäit, wou hie perséinlech dem
glorräiche kaiserleche \VUgras{Marschall} gratuléiere koum, deem säi
\VUgras{Portrait} \textsuperscript{(44)} un der Mauer hänkt an deen den
\VUgras{Urgrousspapp} vu senge Kanner ass.

\VUsaut{6.0mm}
\VUfilet{22mm}
